% Author: SuFan; MIT License. Illustrative dimension samples, not acceptance limits.
\documentclass{article}
\usepackage[margin=20mm]{geometry}
\usepackage{gradbars,booktabs,array}
\usegradbarslibrary{distribution}
\pagestyle{empty}
\begin{document}
% BEGIN TABLE
\begingroup
\gradbarsstyle{batch-demo}{box range=fixed,min=9.8,max=10.2,
  width=70mm,height=5ex,precision=2,label=outside,
  unit={\,mm},palette=accessible}
\newcommand{\batchplot}[1]{\raisebox{\dimexpr.5ex-.5\height+.5\depth\relax}{%
  \gradboxplot[style=batch-demo]{#1}}}
\renewcommand{\arraystretch}{2.2}
\begin{center}
\begin{tabular}{@{}lrp{94mm}@{}}
\toprule
Batch & Samples & 9.80--10.20 mm / median\\
\midrule
A & 9 & \batchplot{9.96,9.98,9.99,10,10.01,10.02,10.03,10.04,10.06}\\
B & 9 & \batchplot{9.85,9.95,9.98,10,10.01,10.02,10.04,10.08,10.17}\\
C & 9 & \batchplot{10.03,10.04,10.05,10.05,10.06,10.07,10.07,10.08,10.09}\\
Pending & -- & \batchplot{NA}\\
\bottomrule
\end{tabular}
\end{center}
\endgroup
% END TABLE
\noindent Every batch uses the same 9.80--10.20 mm display scale. Labels show
medians; hollow outlier marks remain visible. Whiskers are observed values
inside the 1.5 IQR fences, not confidence intervals or acceptance limits.
\end{document}
